\section{Lösungsvorschläge}

Lösungsvorschläge

1. Pins sollten hervorgehoben werden.

Keine der Testpersonen realisierte, dass sie das Menu der Elemente im Downloadtab mit Swipen hervorbringen können. Dies wird gelöst, indem beim erstmaligen öffnen des Downloadtabs ein Element automatisch geswiped wird. So sollten Nutzer mit dieser Funktionalität vertraut werden, ohne dass ihre Nutzung der App beeinträchtigt wird.


Ein Problem für TP4 war, dass Pins nach der Suche nicht hervorgehoben werden. Pins sollten deshalb wenn sie ausgewählt sind hervorgehoben werden, indem ihre Größe verdoppelt wird.

4. Während Suchen getippt werden sollten Suchvorschläge angezeigt werden, so merkt der Nutzer zum einen, dass die Suche funktioniert und zum anderen kann er den gewünschten Pin so schneller finden.

Es war für mehrere Testpersonen problematisch, dass nicht klar war, ob die Suche funktioniert, da es kein Live Feedback gab. Dies wird gelöst, indem Suchvorschläge in beiden Suchfenstern angezeigt werden und indem das Suchmenu im Kartentab automatisch vergrößert wird, während der/die NutzerIn darin tippen. Dies ermöglicht es ihnen auch, weniger zu Tippen.